\section{Principle IX: The Duty of Heritage}

\begin{quotex}
In this particularly beautiful passage, Maurras explains the duty of one generation to pass on its culture, customs, and traditions to the next, which in turn, has the duty to accept and develop that heritage. The loss of one's traditions needs to be taken at least as seriously as the attention we pay today to the possible extinction of various fishes and insects. Unfortunately, the youth of today are the new disinherited, since their fathers have passed on little except a hatred for their own past. How that plays out is uncertain, unless and until they rise up to reclaim what is rightfully theirs. 

\end{quotex}
In order to justify what common sense and custom maintain by force, let us not talk about rights, but instead duty.

The duty to bequeath and make a will.

The duty to inherit.

Those who are called disinherited or proletarian are sometimes understood to be preaching that that is not just in their point of view. That is just and beneficial for everyone.

The good that is established in a natural or legal family can have doubtful sources: it is redeemed and moralized by its fixity, by the firm benefit that it establishes around itself in stabilizing the conditions of life, in distributing work, in securing it, in preparing a point of solid support to the generations to come.

There is no difference between the damages caused to nature by the death of a beautiful animal, followed by the return of its elements to the universal dust, and the destruction of an ordinary fortune to the death of its creator.

There is a dead loss for society as for nature, the obligation to take up again a long and painful task.

Prosperous social organisations are those that prevent these realities from dissolving into the void and which help to keep them from a total death: these purveyors of life, being the born-enemies of destruction, made a respected institution from its heritage, and, one could say, a type of sacrament from its legacy.



\flrightit{Posted on 2012-12-16 by Cologero }

\begin{center}* * *\end{center}

\begin{footnotesize}\begin{sffamily}



\texttt{Jason-Adam on 2012-12-16 at 21:27 said: }

Truly beautiful indeed. The West is in dire straits as even the descendents of the former elite are losing their ancestral heritage…….and instead of rising up to reclaim what is ours, too many faux traditionalists look to another culture….


\end{sffamily}\end{footnotesize}
